\chapter{Revisão da Literatura}
\label{cap:revisao-literatura}

~\citet{survey1} descreve diversas características relacionadas com a detecção, a rastreabilidade e a propagação de \textit{malware} em dispositivos móveis. Além disso, técnicas para interromper a propagação também são discutidas. Com relação à detecção e à rastreabilidade, o protocolo DNS tem papel central, dado que muitos aplicativos maliciosos possuem instruções para acessarem um determinado domínio que contém o \textit{malware}, geralmente oculto em um código aparentemente inofensivo. No geral, é possível rastrear e detectar o \textit{malware} por meio da coleta do tráfego DNS no \textit{gateway} da rede ou utilizando alguma ferramenta de proteção no próprio dispositivo. Nomes de domínios que sejam similares a nomes de organizações conhecidas, que tenham sido recém registrados e que estejam sendo acessados pela primeira vez devem levantar suspeitas e levar a uma pesquisa mais aprofundada. Em casos de confirmação de que o nome de domínio de fato está sendo utilizado para propagação de \textit{malware}, servidores DNS podem ser configurados para realização do chamado DNS \textit{sinkholing}, procedimento em que o endereço IP retornado em uma consulta DNS é um endereço não roteável na Internet, ao invés do endereço real. Dessa forma, o \textit{malware} que esteja armazenado em um site sob um domínio configurado dessa forma, não será copiado, evitando a infecção de dispositivos. Essa técnica tem sido usada com sucesso para interromper ataques, com o mais famoso tendo sido o \textit{ransomware} WannaCry em 2017~\citep{wannacry}. Embora o processo de detecção de nomes de domínios maliciosos esteja bem estabelecido, é importante que isso seja feito de forma automatizada e antes de ataques acontecerem em uma rede em particular, por exemplo por meio de listas de bloqueio que contenham nomes conhecidamente maliciosos, como aquelas disponíveis em \url{https://openphish.com/} e em \url{https://phishtank.org/}. Dessa forma, um ataque que tenha início em um determinado local pode gerar proteções para outros locais.

~\citet{cic-bell-dns-2021} discute o problema de nomes de domínios usados com propósitos maliciosos com ênfase às consequências financeiras causadas por isso. A utilização de listas de bloqueio como possível solução é apresentada, porém essas listas têm o problema de não capturarem domínios recém criados, justificando a sua utilização em combinação com o processo de detecção pelo tráfego da rede de forma descentralizada. Dessa forma, discute-se a utilização de aprendizado de máquina para classificação de tráfego DNS suspeito, com foco nas características dos dados utilizados para treinamento dos modelos. Entre os classificadores avaliados no estudo, que incluem SVM, KNN, MLP, Naive Bayes e Regressão Logística, todos capazes de classificar os domínios em \textit{spam}, \textit{phishing} e \textit{malware}, o algoritmo KNN superou os modelos clássicos de aprendizado de máquina. Os algoritmos que apresentaram melhores desempenhos foram considerados neste TCC. Uma importante contribuição do trabalho é a disponibilização do dataset CIC-Bell-DNS2021, que simula um ambiente real com 99\% de tráfego benigno, utilizado para o treinamento dos classificadores listados e disponível em \url{https://www.unb.ca/cic/datasets/dns-2021.html}. Este projeto de TCC analisou esse dataset, além de outros que foram investigados no decorrer do projeto.

Recentemente, ~\citet{survey2} e ~\citet{survey3} discutiram as diversas soluções apresentadas nos últimos anos para detecção de \textit{malware} e detecção de \textit{phishing} respectivamente. Em ambos os trabalhos, a detecção com base em características do tráfego DNS é listada, com ênfase à limitação na utilização de listas de bloqueio e aos bons resultados que têm sido alcançados com a utilização de modelos de aprendizado de máquina, principalmente aqueles baseados em redes neurais. Observa-se que muitos dos modelos utilizam características dos nomes de domínio que fazem sentido para o idioma Inglês mas não necessariamente para o Português.

~\citet{shafi2024dns} introduz um analisador de fluxo de rede na camada de aplicação, e o novo \textit{dataset} BCCC-CIC-Bell-DNS-2024, que supera as limitações dos datasets DNS públicos existentes até sua data de publicação. Como proposta, este trabalho científico propõe uma abordagem inovadora de perfilamento comportamental de DNS que enfrenta desafios como táticas de evasão, variabilidade de conteúdo, detecção de intenções maliciosas, ofuscação de URLs, ataques de baixa intensidade e a manutenção da precisão diante de comportamentos normais diversos. Além disso, propõe um novo modelo que combina comportamentos únicos de \textit{features} e suas correlações para extração de padrões, seleção de \textit{features} e desenvolvimento de uma arquitetura de rede neural para a precisão na criação de perfis.
Esse dataset, gerado a partir da ferramenta \textit{ALFlowLyzer}~\citep{shafi2024a}, foi o escolhido para conduzir todos os experimentos deste TCC (detalhado no Capítulo~\ref{cap:modelo}), por sua abrangência de mais de 120 \textit{features} e por refletir um panorama mais atual de tráfego DNS malicioso.
Assim como os autores de~\citet{SHAFI2024109436}, a este trabalho é voltado para detecção de tráfego DNS malicioso e utiliza redes neurais. Além disso, foi utilizado o dataset criado. A principal diferença está no fato de que toda a implementação realizada no presente trabalho está disponível como software livre.

~\citet{https://doi.org/10.1002/ett.4150} apresenta uma revisão sistemática da literatura sobre o uso de aprendizado de máquina e aprendizado profundo no desenvolvimento de IDSs. %O estudo avaliou publicações científicas de 2017 a 2020. % a fim de identificar tendências, desafios e direcionamentos futuros. % \\
% Em vista de que um IDS monitora o tráfego de rede para garantir a confidencialidade, integridade e disponibilidade dos dados, o artigo classifica IDSs sob duas perspectivas:\\ - por implantação: HIDS (\textit{Host-based IDS}) e NIDS (\textit{Network-based IDS});\\ - por método de detecção: SIDS (\textit{Signature-based IDS}) e AIDS (\textit{Anomaly-based IDS}).\\
Observou-se que a implementação de um IDS em geral segue três fases fundamentais: pré-processamento de dados, treinamento e teste. Além disso, os autores concluem que há uma tendência na utilização de técnicas de aprendizado profundo baseadas em redes neurais % com bons resultados obtidos,
e que melhorias podem ser alcançadas com % O estudo destaca uma transição tecnológica, na qual 60\% dos novos métodos focam puramente em aprendizado profundo, 20\% em aprendizado tradicional e 20\% em abordagens híbridas. Os principais algoritmos usados são:\\ - aprendizado profundo (DL): Autoencoders, DNN, CNN e RNN;\\ - aprendizado de máquina tradicional (ML): SVM, Random Forest, KNN e DT;\\ - abordagens híbridas (combinação de algoritmos baseados em ML e DL).\\ A eficácia dos modelos é medida através da matriz de confusão, utilizando equações como de: acurácia, precisão, recall e f-measure. Tudo isso é aplicado em conjuntos de dados modernos que refletem melhor as ameaças atuais. As lacunas apontadas pela pesquisa são: o desbalanceamento dos dados (modelos falham em detectar ataques minoritários) e a complexidade de algoritmos e escassez de recursos para modelos que exigem alto poder computacional. E conclui com a afirmação de que
a seleção inteligente de \textit{features}. %, levando à redução da complexidade dos modelos sem sacrificar a precisão da detecção. 
Seguindo as recomendações de~\citet{https://doi.org/10.1002/ett.4150}, este trabalho utiliza redes neurais e seleção inteligente de \textit{features}.

~\citet{sbseg-phishing} apresentam um mecanismo para mitigação de \textit{phishing} baseado na tecnologia eBPF (\textit{extended Berkeley Packet Filter}) e no framework XDP (\textit{eXpress Data Path}). O mecanismo age como um \textit{firewall} de DNS bloqueando resoluções de nomes direcionadas a domínios maliciosos. %Comparando com outras propostas voltadas para o mesmo objetivo, o trabalho consegue bloquear as resoluções com uma redução de cerca de 70\% de consumo de CPU dos servidores e com incremento médio de apenas 12,6$\mu$s na latência das resoluções.
Diferente deste TCC, o trabalho de ~\citet{sbseg-phishing} baseia-se em uma lista de domínios maliciosos pré criada. A detecção de novos domínios não é realizada.